% generated by GAPDoc2LaTeX from XML source (Frank Luebeck)
\documentclass[a4paper,11pt]{report}

            \usepackage{a4wide}
            \newcommand{\bbZ}{\mathbb{Z}}
			\usepackage{mathtools}
        
\usepackage[top=37mm,bottom=37mm,left=27mm,right=27mm]{geometry}
\sloppy
\pagestyle{myheadings}
\usepackage{amssymb}
\usepackage[utf8]{inputenc}
\usepackage{makeidx}
\makeindex
\usepackage{color}
\definecolor{FireBrick}{rgb}{0.5812,0.0074,0.0083}
\definecolor{RoyalBlue}{rgb}{0.0236,0.0894,0.6179}
\definecolor{RoyalGreen}{rgb}{0.0236,0.6179,0.0894}
\definecolor{RoyalRed}{rgb}{0.6179,0.0236,0.0894}
\definecolor{LightBlue}{rgb}{0.8544,0.9511,1.0000}
\definecolor{Black}{rgb}{0.0,0.0,0.0}

\definecolor{linkColor}{rgb}{0.0,0.0,0.554}
\definecolor{citeColor}{rgb}{0.0,0.0,0.554}
\definecolor{fileColor}{rgb}{0.0,0.0,0.554}
\definecolor{urlColor}{rgb}{0.0,0.0,0.554}
\definecolor{promptColor}{rgb}{0.0,0.0,0.589}
\definecolor{brkpromptColor}{rgb}{0.589,0.0,0.0}
\definecolor{gapinputColor}{rgb}{0.589,0.0,0.0}
\definecolor{gapoutputColor}{rgb}{0.0,0.0,0.0}

%%  for a long time these were red and blue by default,
%%  now black, but keep variables to overwrite
\definecolor{FuncColor}{rgb}{0.0,0.0,0.0}
%% strange name because of pdflatex bug:
\definecolor{Chapter }{rgb}{0.0,0.0,0.0}
\definecolor{DarkOlive}{rgb}{0.1047,0.2412,0.0064}


\usepackage{fancyvrb}

\usepackage{mathptmx,helvet}
\usepackage[T1]{fontenc}
\usepackage{textcomp}


\usepackage[
            pdftex=true,
            bookmarks=true,        
            a4paper=true,
            pdftitle={Written with GAPDoc},
            pdfcreator={LaTeX with hyperref package / GAPDoc},
            colorlinks=true,
            backref=page,
            breaklinks=true,
            linkcolor=linkColor,
            citecolor=citeColor,
            filecolor=fileColor,
            urlcolor=urlColor,
            pdfpagemode={UseNone}, 
           ]{hyperref}

\newcommand{\maintitlesize}{\fontsize{50}{55}\selectfont}

% write page numbers to a .pnr log file for online help
\newwrite\pagenrlog
\immediate\openout\pagenrlog =\jobname.pnr
\immediate\write\pagenrlog{PAGENRS := [}
\newcommand{\logpage}[1]{\protect\write\pagenrlog{#1, \thepage,}}
%% were never documented, give conflicts with some additional packages

\newcommand{\GAP}{\textsf{GAP}}

%% nicer description environments, allows long labels
\usepackage{enumitem}
\setdescription{style=nextline}

%% depth of toc
\setcounter{tocdepth}{1}





%% command for ColorPrompt style examples
\newcommand{\gapprompt}[1]{\color{promptColor}{\bfseries #1}}
\newcommand{\gapbrkprompt}[1]{\color{brkpromptColor}{\bfseries #1}}
\newcommand{\gapinput}[1]{\color{gapinputColor}{#1}}


\begin{document}

\logpage{[ 0, 0, 0 ]}
\begin{titlepage}
\mbox{}\vfill

\begin{center}{\maintitlesize \textbf{ LocalActionDiagrams \\
\mbox{}}}\\
\vfill

\hypersetup{pdftitle={ LocalActionDiagrams }}
\markright{\scriptsize \mbox{}\hfill  LocalActionDiagrams  \hfill\mbox{}}
{\Huge \textbf{ Super awesome package for local action diagrams. \\
\mbox{}}}\\
\vfill

{\Huge  0.1 \mbox{}}\\[1cm]
{ 29 January 2023 \mbox{}}\\[1cm]
\mbox{}\\[2cm]
{\Large \textbf{\strut  Marcus Chijoff     \strut\mbox{}}}\\
{\Large \textbf{\strut  Stephan Tornier     \strut\mbox{}}}\\
\hypersetup{pdfauthor={ Marcus Chijoff    ;  Stephan Tornier    }}
\end{center}\vfill

\mbox{}\\
{\mbox{}\\
\small \noindent \textbf{ Marcus Chijoff    }  Email: \href{mailto://marcus.chijoff@uon.edu.au} {\texttt{marcus.chijoff@uon.edu.au}}\\
  Homepage: \href{https://newcastle.edu.au} {\texttt{https://newcastle.edu.au}}\\
  Address: \begin{minipage}[t]{8cm}\noindent
 No\\
 \end{minipage}
}\\
{\mbox{}\\
\small \noindent \textbf{ Stephan Tornier    }  Email: \href{mailto://stephan.tornier@newcastle.edu.au} {\texttt{stephan.tornier@newcastle.edu.au}}\\
  Homepage: \href{https://newcastle.edu.au} {\texttt{https://newcastle.edu.au}}\\
  Address: \begin{minipage}[t]{8cm}\noindent
 No\\
 \end{minipage}
}\\
\end{titlepage}

\newpage\setcounter{page}{2}
\newpage

\def\contentsname{Contents\logpage{[ 0, 0, 1 ]}}

\tableofcontents
\newpage

     
\chapter{\textcolor{Chapter }{RSGraphs}}\label{Chapter_RSGraphs}
\logpage{[ 1, 0, 0 ]}
\hyperdef{L}{X7AC2E3EF7B93A9FA}{}
{
  

 An RSGraph is a graph $\Gamma = (V, A, o, r)$ that consists of a vertex set $V$, an arc set $A$, a map $o : A \to V$ which assigns each vertex to an \emph{origin}, and a bijection $r : A \to A$ which assigns each arc to a \emph{reverse} arc. There is also a map $t := o \circ r$ which assigns each arc to a \emph{terminus} vertex. Note the RSGraphs are allowed to contain loops and parallel arcs. 

 We created the RSGraph object for this package due to our use of graphs having
a specialised definition. For example, every arc in an RSGraph must have a
reverse arc associated with it while a \emph{Digraph} from the \textsf{Digraphs} allows for arcs that do not have an associated reverse arc. Creating a new
object category for RSGraphs allows for both better error checking when using
RSGraphs and to avoid conflicting definitions, such as the definition of a
cycle for a Digraph and an RSGraph. We provide functionality to convert
between an RSGraph and a \emph{Graph} from the \textsf{GRAPE} package or a \emph{Digraph} from the \textsf{Digraphs} package. 

 The name "RSGraph" arises from Reid and Smith, the authors of the local action
diagram paper. This name was chosen to avoid conflicting with the \textsf{GRAPE} package. 

 
\section{\textcolor{Chapter }{Creating RSGraphs}}\label{Chapter_RSGraphs_Section_Creating_RSGraphs}
\logpage{[ 1, 1, 0 ]}
\hyperdef{L}{X7AA329AC8054CB4D}{}
{
  

 

\subsection{\textcolor{Chapter }{IsRSGraph (for an object)}}
\logpage{[ 1, 1, 1 ]}\nobreak
\hyperdef{L}{X7D238B008758C147}{}
{\noindent\textcolor{FuncColor}{$\triangleright$\enspace\texttt{IsRSGraph({\mdseries\slshape obj})\index{IsRSGraph@\texttt{IsRSGraph}!for an object}
\label{IsRSGraph:for an object}
}\hfill{\scriptsize (Category)}}\\
\textbf{\indent Returns:\ }
\texttt{true} if \mbox{\texttt{\mdseries\slshape graph}} is of the category \texttt{IsRSGraph} and \texttt{false} otherwise. 



 Every RSGraph object belongs to the category \texttt{IsRSGraph}. Furthermore, every RSGraph is an immutable attribute storing object. }

 

\subsection{\textcolor{Chapter }{RSGraphByAdjacencyList (for a list, permutation[, and list])}}
\logpage{[ 1, 1, 2 ]}\nobreak
\label{AdjacencyList}
\hyperdef{L}{X794620DB807A22B3}{}
{\noindent\textcolor{FuncColor}{$\triangleright$\enspace\texttt{RSGraphByAdjacencyList({\mdseries\slshape adjacency{\textunderscore}list, reverse{\textunderscore}map[, vertex{\textunderscore}ids]})\index{RSGraphByAdjacencyList@\texttt{RSGraphByAdjacencyList}!for a list, permutation[, and list]}
\label{RSGraphByAdjacencyList:for a list, permutation[, and list]}
}\hfill{\scriptsize (operation)}}\\
\noindent\textcolor{FuncColor}{$\triangleright$\enspace\texttt{RSGraphByAdjacencyListNC({\mdseries\slshape adjacency{\textunderscore}list, reverse{\textunderscore}map[, vertex{\textunderscore}ids]})\index{RSGraphByAdjacencyListNC@\texttt{RSGraphByAdjacencyListNC}!for a list, permutation[, and list]}
\label{RSGraphByAdjacencyListNC:for a list, permutation[, and list]}
}\hfill{\scriptsize (operation)}}\\
\textbf{\indent Returns:\ }
An RSGraph. 



 This function creates an RSGraph described by an adjacency list. The argument \mbox{\texttt{\mdseries\slshape adjacnecy{\textunderscore}list}} is a list of the from \texttt{[[origin, terminus], ...]}. Each arc in the graph is given an id which corresponds to its position in
the list \mbox{\texttt{\mdseries\slshape adjacency{\textunderscore}list}}. 

 The reverse map must be a permutation that is an involution
\texttt{\symbol{45}}\texttt{\symbol{45}}\texttt{\symbol{45}} i.e. \texttt{reverse{\textunderscore}map*reverse{\textunderscore}map = ()}. Furthermore, if element \mbox{\texttt{\mdseries\slshape idx}} of \mbox{\texttt{\mdseries\slshape adjacency{\textunderscore}list}} is \texttt{[u, v]} then \mbox{\texttt{\mdseries\slshape adjacency{\textunderscore}list[idx\texttt{\symbol{94}}reverse{\textunderscore}map]}} must equal \texttt{[v, u]}. 

 By default, the vertices are labelled with ids $\{1, 2, \dots, n\}$ where $n$ is the maximal element of \texttt{Flat(adjacency{\textunderscore}list)}. If the optional third argument \mbox{\texttt{\mdseries\slshape vertex{\textunderscore}ids}} is given then these are taken to be the vertex ids. This argument must be a
dense list of integers. 

 Finally, the underlying graph must be connected. 

 The NC variant of the function does not check that the graph is connected, the
reverse map is valid, and that each arc has an associated reverse arc. }

 
\begin{Verbatim}[commandchars=!@|,fontsize=\small,frame=single,label=Example]
  !gapprompt@gap>| !gapinput@adj_list := [[1, 2], [2, 1], [2, 2], [2, 2], [1, 1]];;|
  !gapprompt@gap>| !gapinput@rev_map := (1,2)(3,4);;|
  !gapprompt@gap>| !gapinput@graph := RSGraphByAdjacencyList(adj_list, rev_map);|
  <RSGraph with 2 vertices and 5 arcs>
  
  !gapprompt@gap>| !gapinput@adj_list2 := [[2, 2]];;|
  !gapprompt@gap>| !gapinput@rev_map2 := ();;|
  !gapprompt@gap>| !gapinput@vertex_ids := [2];;|
  !gapprompt@gap>| !gapinput@graph := RSGraphByAdjacencyList(adj_list, rev_map, vertex_ids);|
  <RSGraph with 1 vertex and 1 arc>
\end{Verbatim}
 

\subsection{\textcolor{Chapter }{RSGraphByAdjacencyMatrix (for a matrix, permutation[, and list])}}
\logpage{[ 1, 1, 3 ]}\nobreak
\label{AdjacencyMatrix}
\hyperdef{L}{X873AEEA579B2330C}{}
{\noindent\textcolor{FuncColor}{$\triangleright$\enspace\texttt{RSGraphByAdjacencyMatrix({\mdseries\slshape adjacency{\textunderscore}matrix, reverse{\textunderscore}map[, vertex{\textunderscore}ids]})\index{RSGraphByAdjacencyMatrix@\texttt{RSGraphByAdjacencyMatrix}!for a matrix, permutation[, and list]}
\label{RSGraphByAdjacencyMatrix:for a matrix, permutation[, and list]}
}\hfill{\scriptsize (operation)}}\\
\noindent\textcolor{FuncColor}{$\triangleright$\enspace\texttt{RSGraphByAdjacencyMatrixNC({\mdseries\slshape adjacency{\textunderscore}matrix, reverse{\textunderscore}map[, vertex{\textunderscore}ids]})\index{RSGraphByAdjacencyMatrixNC@\texttt{RSGraphByAdjacencyMatrixNC}!for a matrix, permutation[, and list]}
\label{RSGraphByAdjacencyMatrixNC:for a matrix, permutation[, and list]}
}\hfill{\scriptsize (operation)}}\\
\textbf{\indent Returns:\ }
An RSGraph. 



 This function creates an RSGraph described by an adjacency matrix. The
argument \mbox{\texttt{\mdseries\slshape adjacnecy{\textunderscore}matrix}} is an $n \times n$ matrix represented by a list of lists. This means that there are $m$ arcs between vertices $u$ and $v$ of the graph if \texttt{adjacency{\textunderscore}matrix[u][v] = m}. 

 The matrix is read in row\texttt{\symbol{45}}major order to determine the arc
ids of the graph. This is important for determining the reverse map of the
graph. The reverse map must be a permutation that is an involution
\texttt{\symbol{45}}\texttt{\symbol{45}}\texttt{\symbol{45}} i.e. \texttt{reverse{\textunderscore}map*reverse{\textunderscore}map = ()}. 

 By default, the vertices are labelled with ids $\{1, 2, \dots, n\}$ where $n$ \texttt{Size(adjacency{\textunderscore}matrix)} \texttt{\symbol{45}}\texttt{\symbol{45}}\texttt{\symbol{45}} i.e. the number
of rows (and columns) of the matrix. If the optional third argument \mbox{\texttt{\mdseries\slshape vertex{\textunderscore}ids}} is given then these are taken to be the vertex ids. This argument must be a
dense list of integers and element \texttt{idx} of the list will be the label of the vertex in row and column \texttt{idx} of \mbox{\texttt{\mdseries\slshape adjacency{\textunderscore}matrix}}. 

 Finally, the underlying graph must be connected. 

 Note that the first argument must be a list of lists with each sublist
containing integers. In particular, it can not be an object in the category \texttt{IsMatrixObj} constructed by the function \texttt{Matrix}. 

 The NC variant of the function does not check that the graph is connected, the
reverse map is valid, and that each arc has an associated reverse arc. }

 
\begin{Verbatim}[commandchars=!@|,fontsize=\small,frame=single,label=Example]
  !gapprompt@gap>| !gapinput@adj_mat := [[1, 2], [2, 0]];;|
  !gapprompt@gap>| !gapinput@rev_map := (2, 4)(3,5);;|
  !gapprompt@gap>| !gapinput@vertex_ids := [2,3];;|
  !gapprompt@gap>| !gapinput@graph := RSGraphByAdjacencyMatrix(adj_list, rev_map, vertex_ids);|
  <RSGraph with 2 vertices and 5 arcs>
  
  !gapprompt@gap>| !gapinput@Print(graph);|
  Vertices = { 2, 3 }
  Arcs = {
          1 = ( origin = 2, terminus = 2, inverse = 1 )
          2 = ( origin = 2, terminus = 3, inverse = 4 )
          3 = ( origin = 2, terminus = 3, inverse = 5 )
          4 = ( origin = 3, terminus = 2, inverse = 2 )
          5 = ( origin = 3, terminus = 2, inverse = 3 )
  }
  Reverse Map = (2,4)(3,5)
\end{Verbatim}
 }

 
\section{\textcolor{Chapter }{RSGraph Attributes and Properties}}\label{Chapter_RSGraphs_Section_RSGraph_Attributes_and_Properties}
\logpage{[ 1, 2, 0 ]}
\hyperdef{L}{X7AED02578003048A}{}
{
  

 

\subsection{\textcolor{Chapter }{RSGraphVertices}}
\logpage{[ 1, 2, 1 ]}\nobreak
\hyperdef{L}{X81642D0C7C970254}{}
{\noindent\textcolor{FuncColor}{$\triangleright$\enspace\texttt{RSGraphVertices({\mdseries\slshape graph})\index{RSGraphVertices@\texttt{RSGraphVertices}}
\label{RSGraphVertices}
}\hfill{\scriptsize (attribute)}}\\
\textbf{\indent Returns:\ }
The list of vertex ids of \mbox{\texttt{\mdseries\slshape graph}}. 



 

 }

 

\subsection{\textcolor{Chapter }{RSGraphNumberVertices}}
\logpage{[ 1, 2, 2 ]}\nobreak
\hyperdef{L}{X7DF76A7B81FE962B}{}
{\noindent\textcolor{FuncColor}{$\triangleright$\enspace\texttt{RSGraphNumberVertices({\mdseries\slshape graph})\index{RSGraphNumberVertices@\texttt{RSGraphNumberVertices}}
\label{RSGraphNumberVertices}
}\hfill{\scriptsize (attribute)}}\\
\textbf{\indent Returns:\ }
The number of vertices of \mbox{\texttt{\mdseries\slshape graph}}. 



 

 }

 

\subsection{\textcolor{Chapter }{RSGraphArcs}}
\logpage{[ 1, 2, 3 ]}\nobreak
\hyperdef{L}{X80C7D0FB80332FCE}{}
{\noindent\textcolor{FuncColor}{$\triangleright$\enspace\texttt{RSGraphArcs({\mdseries\slshape graph})\index{RSGraphArcs@\texttt{RSGraphArcs}}
\label{RSGraphArcs}
}\hfill{\scriptsize (attribute)}}\\
\textbf{\indent Returns:\ }
The record of arcs of \mbox{\texttt{\mdseries\slshape graph}}. 



 The name of each component of the record is the id of the arc. Each component
it itself a record with three components: origin, terminus, and inverse. These
store the origin vertex, terminus vertex, and inverse arc id respectively. 

 For iterating over the arcs of an RSGraph see \texttt{RSGraphArcIterator} (\ref{RSGraphArcIterator}). }

 
\begin{Verbatim}[commandchars=!@|,fontsize=\small,frame=single,label=Example]
  !gapprompt@gap>| !gapinput@adj_list := [[1, 2], [2, 1], [2, 2], [2, 2], [1, 1]];;|
  !gapprompt@gap>| !gapinput@rev_map := (1,2)(3,4);;|
  !gapprompt@gap>| !gapinput@graph := RSGraphByAdjacencyList(adj_list, rev_map);|
  <RSGraph with 2 vertices and 5 arcs>
  !gapprompt@gap>| !gapinput@RSGraphArcs(graph).1;|
  rec( inverse := 2, origin := 1, terminus := 2 )
\end{Verbatim}
 

\subsection{\textcolor{Chapter }{RSGraphNumberArcs}}
\logpage{[ 1, 2, 4 ]}\nobreak
\hyperdef{L}{X7FB8ECE0858F72D2}{}
{\noindent\textcolor{FuncColor}{$\triangleright$\enspace\texttt{RSGraphNumberArcs({\mdseries\slshape graph})\index{RSGraphNumberArcs@\texttt{RSGraphNumberArcs}}
\label{RSGraphNumberArcs}
}\hfill{\scriptsize (attribute)}}\\
\textbf{\indent Returns:\ }
The number of arcs of \mbox{\texttt{\mdseries\slshape graph}}. 



 

 }

 

\subsection{\textcolor{Chapter }{RSGraphArcIDs}}
\logpage{[ 1, 2, 5 ]}\nobreak
\hyperdef{L}{X833D53A5856F0FCE}{}
{\noindent\textcolor{FuncColor}{$\triangleright$\enspace\texttt{RSGraphArcIDs({\mdseries\slshape graph})\index{RSGraphArcIDs@\texttt{RSGraphArcIDs}}
\label{RSGraphArcIDs}
}\hfill{\scriptsize (attribute)}}\\
\textbf{\indent Returns:\ }
The list of arc ids of \mbox{\texttt{\mdseries\slshape graph}}. 



 

 }

 

\subsection{\textcolor{Chapter }{RSGraphReverseMap}}
\logpage{[ 1, 2, 6 ]}\nobreak
\hyperdef{L}{X79E3EA558513C547}{}
{\noindent\textcolor{FuncColor}{$\triangleright$\enspace\texttt{RSGraphReverseMap({\mdseries\slshape graph})\index{RSGraphReverseMap@\texttt{RSGraphReverseMap}}
\label{RSGraphReverseMap}
}\hfill{\scriptsize (attribute)}}\\
\textbf{\indent Returns:\ }
The reverse map of \mbox{\texttt{\mdseries\slshape graph}}. 



 

 }

 

\subsection{\textcolor{Chapter }{RSGraphAdjacencyMatrix}}
\logpage{[ 1, 2, 7 ]}\nobreak
\hyperdef{L}{X8574F86678FB772F}{}
{\noindent\textcolor{FuncColor}{$\triangleright$\enspace\texttt{RSGraphAdjacencyMatrix({\mdseries\slshape graph})\index{RSGraphAdjacencyMatrix@\texttt{RSGraphAdjacencyMatrix}}
\label{RSGraphAdjacencyMatrix}
}\hfill{\scriptsize (attribute)}}\\
\textbf{\indent Returns:\ }
The adjacency matrix of \mbox{\texttt{\mdseries\slshape graph}}. 



 

 }

 

\subsection{\textcolor{Chapter }{RSGraphOutNeighbours}}
\logpage{[ 1, 2, 8 ]}\nobreak
\hyperdef{L}{X80D76940852FBF6D}{}
{\noindent\textcolor{FuncColor}{$\triangleright$\enspace\texttt{RSGraphOutNeighbours({\mdseries\slshape graph})\index{RSGraphOutNeighbours@\texttt{RSGraphOutNeighbours}}
\label{RSGraphOutNeighbours}
}\hfill{\scriptsize (attribute)}}\\
\textbf{\indent Returns:\ }
The record of vertex neighbours of \mbox{\texttt{\mdseries\slshape graph}}. 



 The name of each component is a vertex id of the graph. Each element is a list
of vertex ids corresponding to the neighbours of the vertex. A vertex id \texttt{v{\textunderscore}id} is in list \texttt{idx} if and only if there is an arc connecting the vertices \texttt{idx} and \texttt{v{\textunderscore}id}. For example, if there is an arc from vertex 1 to 3 of the graph then 3 is an
element of \texttt{RSGraphOutNeighbours(graph).1}. }

 

\subsection{\textcolor{Chapter }{RSGraphInNeighbours}}
\logpage{[ 1, 2, 9 ]}\nobreak
\hyperdef{L}{X7AC26C127F7A4752}{}
{\noindent\textcolor{FuncColor}{$\triangleright$\enspace\texttt{RSGraphInNeighbours({\mdseries\slshape graph})\index{RSGraphInNeighbours@\texttt{RSGraphInNeighbours}}
\label{RSGraphInNeighbours}
}\hfill{\scriptsize (attribute)}}\\
\textbf{\indent Returns:\ }
The record of vertex neighbours of \mbox{\texttt{\mdseries\slshape graph}}. 



 Since every arc has an inverse arc this returns the same result as \texttt{RSGraphOutNeighbours(graph)}. }

 

\subsection{\textcolor{Chapter }{RSGraphOutArcs}}
\logpage{[ 1, 2, 10 ]}\nobreak
\hyperdef{L}{X7D3AFF9E780338C7}{}
{\noindent\textcolor{FuncColor}{$\triangleright$\enspace\texttt{RSGraphOutArcs({\mdseries\slshape graph})\index{RSGraphOutArcs@\texttt{RSGraphOutArcs}}
\label{RSGraphOutArcs}
}\hfill{\scriptsize (attribute)}}\\
\textbf{\indent Returns:\ }
The record of arcs originating at each vertex. 



 The name of each component is a vertex id of the graph. Each element is a list
of arc ids. An arc id is in list \texttt{idx} if and only if the origin of that arc is \texttt{idx}. For example, if an arc with id \texttt{2} originates at vertex 3 then \texttt{2} is an element of \texttt{RSGraphInNeighbours(graph).3}. }

 

\subsection{\textcolor{Chapter }{RSGraphInArcs}}
\logpage{[ 1, 2, 11 ]}\nobreak
\hyperdef{L}{X78FADD068325B2C0}{}
{\noindent\textcolor{FuncColor}{$\triangleright$\enspace\texttt{RSGraphInArcs({\mdseries\slshape graph})\index{RSGraphInArcs@\texttt{RSGraphInArcs}}
\label{RSGraphInArcs}
}\hfill{\scriptsize (attribute)}}\\
\textbf{\indent Returns:\ }
The record of arcs originating at each vertex. 



 The name of each component is a vertex id of the graph. Each element is a list
of arc ids. An arc id is in list \texttt{idx} if and only if the terminus of that arc is \texttt{idx}. For example, if an arc with id \texttt{2} terminates at vertex 1 then \texttt{2} is an element of \texttt{RSGraphInNeighbours(graph).1}. }

 
\begin{Verbatim}[commandchars=!@|,fontsize=\small,frame=single,label=Example]
  !gapprompt@gap>| !gapinput@adj_list := [[1, 2], [2, 1], [2, 2], [2, 2]];;|
  !gapprompt@gap>| !gapinput@rev_map := (1,2)(3,4);;|
  !gapprompt@gap>| !gapinput@graph := RSGraphByAdjacencyList(adj_list, rev_map);|
  <RSGraph with 2 vertices and 4 arcs>
  !gapprompt@gap>| !gapinput@RSGraphOutNeighbours(graph);|
  rec( 1 := [ 2 ], 2 := [ 1, 2 ] )
  !gapprompt@gap>| !gapinput@RSGraphInNeighbours(graph);|
  rec( 1 := [ 2 ], 2 := [ 1, 2 ] )
  !gapprompt@gap>| !gapinput@RSGraphOutArcs(graph);|
  rec( 1 := [ 1 ], 2 := [ 2, 3, 4 ] )
  !gapprompt@gap>| !gapinput@RSGraphInArcs(graph);|
  rec( 1 := [ 2 ], 2 := [ 1, 3, 4 ] )
\end{Verbatim}
 

\subsection{\textcolor{Chapter }{RSGraphBipartition}}
\logpage{[ 1, 2, 12 ]}\nobreak
\hyperdef{L}{X7B3A52EA86D782F8}{}
{\noindent\textcolor{FuncColor}{$\triangleright$\enspace\texttt{RSGraphBipartition({\mdseries\slshape graph})\index{RSGraphBipartition@\texttt{RSGraphBipartition}}
\label{RSGraphBipartition}
}\hfill{\scriptsize (attribute)}}\\
\textbf{\indent Returns:\ }
The bipartition of \texttt{graph} if it is bipartite and \texttt{fail} otherwise. 



 If the graph is bipartite then this returns a list containing two lists. These
two lists contain the vertex ids in each bipartition. 

 Note that for an RSGraph to be bipartite it must not contain parallel edges or
loops. }

 
\begin{Verbatim}[commandchars=!@|,fontsize=\small,frame=single,label=Example]
  !gapprompt@gap>| !gapinput@adj_list := [[1, 2], [2, 1], [1, 3], [3, 1]];;|
  !gapprompt@gap>| !gapinput@rev_map := (1,2)(3,4);;|
  !gapprompt@gap>| !gapinput@graph := RSGraphByAdjacencyList(adj_list, rev_map);|
  <RSGraph with 3 vertices and 4 arcs>
  !gapprompt@gap>| !gapinput@RSGraphBipartition(graph);|
  [ [ 1 ], [ 2, 3 ] ]
  
  !gapprompt@gap>| !gapinput@adj_list := [[1, 2], [2, 1], [2, 2], [2, 2]];;|
  !gapprompt@gap>| !gapinput@rev_map := (1,2)(3,4);;|
  !gapprompt@gap>| !gapinput@graph := RSGraphByAdjacencyList(adj_list, rev_map);|
  <RSGraph with 2 vertices and 4 arcs>
  !gapprompt@gap>| !gapinput@RSGraphBipartition(graph);|
  fail
\end{Verbatim}
 

\subsection{\textcolor{Chapter }{RSGraphDegree}}
\logpage{[ 1, 2, 13 ]}\nobreak
\hyperdef{L}{X7C306BFA7DDCA0BA}{}
{\noindent\textcolor{FuncColor}{$\triangleright$\enspace\texttt{RSGraphDegree({\mdseries\slshape graph})\index{RSGraphDegree@\texttt{RSGraphDegree}}
\label{RSGraphDegree}
}\hfill{\scriptsize (attribute)}}\\
\textbf{\indent Returns:\ }
The maximum number of arcs originating/terminating at any vertex in the graph. 



 

 }

 

\subsection{\textcolor{Chapter }{RSGraphIsCycle}}
\logpage{[ 1, 2, 14 ]}\nobreak
\hyperdef{L}{X84FE37BF7CABC401}{}
{\noindent\textcolor{FuncColor}{$\triangleright$\enspace\texttt{RSGraphIsCycle({\mdseries\slshape graph})\index{RSGraphIsCycle@\texttt{RSGraphIsCycle}}
\label{RSGraphIsCycle}
}\hfill{\scriptsize (property)}}\\
\textbf{\indent Returns:\ }
\texttt{true} if \mbox{\texttt{\mdseries\slshape graph}} is a cycle graph and \texttt{false} otherwise. 



 This is a cycle in the sense of !!!cite RS paper here!!!. A graph is a cycle
if: 
\begin{itemize}
\item  It has a single vertex, two arcs, and a non\texttt{\symbol{45}}trivial reverse
map. 
\item  It has two vertices and two edges (i.e. four arcs) between them. 
\item  It has three or more vertices, no loops, and all vertices have degree two. 
\end{itemize}
 }

 

\subsection{\textcolor{Chapter }{RSGraphIsBipartite}}
\logpage{[ 1, 2, 15 ]}\nobreak
\hyperdef{L}{X81C8BFF07A3F1397}{}
{\noindent\textcolor{FuncColor}{$\triangleright$\enspace\texttt{RSGraphIsBipartite({\mdseries\slshape graph})\index{RSGraphIsBipartite@\texttt{RSGraphIsBipartite}}
\label{RSGraphIsBipartite}
}\hfill{\scriptsize (property)}}\\
\textbf{\indent Returns:\ }
\texttt{true} if \mbox{\texttt{\mdseries\slshape graph}} is a bipartite graph and \texttt{false} otherwise. 



 This is equivalent to \texttt{RSGraphBipartition(graph) {\textless}{\textgreater} fail}. }

 

\subsection{\textcolor{Chapter }{RSGraphHasParallelArcs}}
\logpage{[ 1, 2, 16 ]}\nobreak
\hyperdef{L}{X7B1E96FF7ACFF82C}{}
{\noindent\textcolor{FuncColor}{$\triangleright$\enspace\texttt{RSGraphHasParallelArcs({\mdseries\slshape graph})\index{RSGraphHasParallelArcs@\texttt{RSGraphHasParallelArcs}}
\label{RSGraphHasParallelArcs}
}\hfill{\scriptsize (property)}}\\
\textbf{\indent Returns:\ }
\texttt{true} if \mbox{\texttt{\mdseries\slshape graph}} has multiple arcs in the same direction between any two pairs of vertices. 



 

 }

 }

 
\section{\textcolor{Chapter }{RSGraph Operations}}\label{Chapter_RSGraphs_Section_RSGraph_Operations}
\logpage{[ 1, 3, 0 ]}
\hyperdef{L}{X7A35D9E480C3065F}{}
{
  

 

\subsection{\textcolor{Chapter }{RSGraphArcIterator}}
\logpage{[ 1, 3, 1 ]}\nobreak
\hyperdef{L}{X868F3DEE7BECC8A4}{}
{\noindent\textcolor{FuncColor}{$\triangleright$\enspace\texttt{RSGraphArcIterator({\mdseries\slshape graph})\index{RSGraphArcIterator@\texttt{RSGraphArcIterator}}
\label{RSGraphArcIterator}
}\hfill{\scriptsize (operation)}}\\
\textbf{\indent Returns:\ }
An iterator which iterates over the arcs of \mbox{\texttt{\mdseries\slshape graph}}. 



 This functions returns an iterator object which iterates over the arcs of \mbox{\texttt{\mdseries\slshape graph}}. Each element of the iterator is of the form \texttt{[arc{\textunderscore}id, arc{\textunderscore}record]} where \texttt{arc{\textunderscore}record} is in the form described in \texttt{RSGraphArcs} (\ref{RSGraphArcs}). }

 
\begin{Verbatim}[commandchars=!@|,fontsize=\small,frame=single,label=Example]
  !gapprompt@gap>| !gapinput@adj_list := [[1, 2], [2, 1], [2, 2]];;|
  !gapprompt@gap>| !gapinput@rev_map := (1,2);;|
  !gapprompt@gap>| !gapinput@graph := RSGraphByAdjacencyList(adj_list, rev_map);;|
  !gapprompt@gap>| !gapinput@for arc in RSGraphArcIterator(graph) do|
  !gapprompt@gap>| !gapinput@    View(arc); |
  !gapprompt@gap>| !gapinput@    Print("\n");|
  !gapprompt@gap>| !gapinput@od; |
  [ 1, rec( inverse := 2, origin := 1, terminus := 2 ) ]
  [ 2, rec( inverse := 1, origin := 2, terminus := 1 ) ]
  [ 3, rec( inverse := 3, origin := 2, terminus := 2 ) ]
\end{Verbatim}
 

 

\subsection{\textcolor{Chapter }{RSGraphSubgraph (for an RSGraph and List of arcs)}}
\logpage{[ 1, 3, 2 ]}\nobreak
\label{Subgraph}
\hyperdef{L}{X7B622C9F7B37448B}{}
{\noindent\textcolor{FuncColor}{$\triangleright$\enspace\texttt{RSGraphSubgraph({\mdseries\slshape graph, arc{\textunderscore}list})\index{RSGraphSubgraph@\texttt{RSGraphSubgraph}!for an RSGraph and List of arcs}
\label{RSGraphSubgraph:for an RSGraph and List of arcs}
}\hfill{\scriptsize (operation)}}\\
\noindent\textcolor{FuncColor}{$\triangleright$\enspace\texttt{RSGraphSubgraphNC({\mdseries\slshape graph, arc{\textunderscore}list})\index{RSGraphSubgraphNC@\texttt{RSGraphSubgraphNC}!for an RSGraph and List of arcs}
\label{RSGraphSubgraphNC:for an RSGraph and List of arcs}
}\hfill{\scriptsize (operation)}}\\
\textbf{\indent Returns:\ }
An arc induced subgraph of \mbox{\texttt{\mdseries\slshape graph}}. 



 Given a graph and list of arcs in this graph, this function returns the
subgraph consisting exactly of those arcs and the vertices connecting them.
The vertex ids, arc ids, and reverse map are preserved. Note that the reverse
map is not changed and so may act on a larger set of integers than is strictly
necessary. 

 If an arc is in \mbox{\texttt{\mdseries\slshape arc{\textunderscore}list}} then its reverse must also be in it. Furthermore, the resulting subgraph must
be connected. 

 The NC variant of this function does not check for connectivity or for the
inclusion of arc reversals. }

 
\begin{Verbatim}[commandchars=!@|,fontsize=\small,frame=single,label=Example]
  !gapprompt@gap>| !gapinput@adj_list := [[1, 2], [2, 1], [1, 2], [2, 1], [2, 2]];;|
  !gapprompt@gap>| !gapinput@rev_map := (1,2)(3,4);;|
  !gapprompt@gap>| !gapinput@graph := RSGraphByAdjacencyList(adj_list, rev_map);;|
  !gapprompt@gap>| !gapinput@Print(graph);|
  Vertices = { 1, 2 }
  Arcs = {
          1 = ( origin = 1, terminus = 2, inverse = 2 )
          2 = ( origin = 2, terminus = 1, inverse = 1 )
          3 = ( origin = 1, terminus = 2, inverse = 4 )
          4 = ( origin = 2, terminus = 1, inverse = 3 )
          5 = ( origin = 2, terminus = 2, inverse = 5 )
  }
  Reverse Map = (1,2)(3,4)
  !gapprompt@gap>| !gapinput@subgraph := RSGraphSubgraph(graph, [1,2,5]);|
  Print(subgraph);
  Vertices = { 1, 2 }
  Arcs = {
          1 = ( origin = 1, terminus = 2, inverse = 2 )
          2 = ( origin = 2, terminus = 1, inverse = 1 )
          5 = ( origin = 2, terminus = 2, inverse = 5 )
  }
  Reverse Map = (1,2)(3,4)
\end{Verbatim}
 

 

\subsection{\textcolor{Chapter }{RSGraphToStandardForm}}
\logpage{[ 1, 3, 3 ]}\nobreak
\hyperdef{L}{X8011A75C80C0DBF0}{}
{\noindent\textcolor{FuncColor}{$\triangleright$\enspace\texttt{RSGraphToStandardForm({\mdseries\slshape graph})\index{RSGraphToStandardForm@\texttt{RSGraphToStandardForm}}
\label{RSGraphToStandardForm}
}\hfill{\scriptsize (operation)}}\\
\textbf{\indent Returns:\ }
A record containing the graph in a standard range and the maps from the old
graph to the new graph. 



 Given an arbitrary RSGraph with $N$ vertices and $M$ arcs this function changes the vertex ids to be in $\{1, 2, \dots, N\}$ and arc ids to be in the range $\{1, 2, \dots, M\}$. It also ensures that the arcs are sorted in lexicographical order with
respect to the origin and terminus vertices. 

 If $\Gamma$ is the \mbox{\texttt{\mdseries\slshape graph}} then the vertex mapping is defined by $V(\Gamma)_i \mapsto i$. The arcs origin and terminus vertices then have this mapping applied to
them. The resulting arcs are then sorted in lexicographical order so that an
arc with a smaller arc id will have an origin vertex with a smaller id than
one with a larger arc id. If the origin vertices are equal then the arc with
the smaller terminus id will have a smaller arc id. The reverse map of the
graph is changed to have the same action on the new arc ids. 

 This function returns a record containing three components. The \texttt{graph} component contains the new graph constructed from it, the \texttt{vertex{\textunderscore}id{\textunderscore}map} component contains the map from the vertex ids of the original graph to the
new graph, and the \texttt{arc{\textunderscore}id{\textunderscore}map} component contains the map from the arc ids of the original graph to the new
graph. 

 The \texttt{vertex{\textunderscore}id{\textunderscore}map} and \texttt{arc{\textunderscore}id{\textunderscore}map} are both in the category \texttt{IsConstantTimeAccessGeneralMapping}. }

 
\begin{Verbatim}[commandchars=!@|,fontsize=\small,frame=single,label=Example]
  !gapprompt@gap>| !gapinput@adj_list := [[2, 3], [3, 2], [2, 3], [3, 2], [3, 3]];;|
  !gapprompt@gap>| !gapinput@rev_map := (1,2)(3,4);;|
  !gapprompt@gap>| !gapinput@verrtex_ids := [2, 3];;|
  !gapprompt@gap>| !gapinput@graph := RSGraphByAdjacencyList(adj_list, rev_map, vertex_ids);;|
  !gapprompt@gap>| !gapinput@Print(graph);|
  Vertices = { 2, 3 }
  Arcs = {
          1 = ( origin = 2, terminus = 3, inverse = 2 )
          2 = ( origin = 3, terminus = 2, inverse = 1 )
          3 = ( origin = 2, terminus = 3, inverse = 4 )
          4 = ( origin = 3, terminus = 2, inverse = 3 )
          5 = ( origin = 3, terminus = 2, inverse = 5 )
  }
  Reverse Map = (1,2)(3,4)
  !gapprompt@gap>| !gapinput@standard_rec := RSGraphToStandardForm(graph);;|
  !gapprompt@gap>| !gapinput@Print(standard_rec.graph);|
  Vertices = { 1, 2 }
  Arcs = {
          1 = ( origin = 1, terminus = 2, inverse = 3 )
          2 = ( origin = 1, terminus = 2, inverse = 4 )
          3 = ( origin = 2, terminus = 1, inverse = 1 )
          4 = ( origin = 2, terminus = 1, inverse = 2 )
          5 = ( origin = 2, terminus = 2, inverse = 5 )
  }
  Reverse Map = (1,3)(2,4)
  !gapprompt@gap>| !gapinput@Print(2^standard_rec.arc_id_map);|
  3
\end{Verbatim}
 

\subsection{\textcolor{Chapter }{RSGraphSpanningTree}}
\logpage{[ 1, 3, 4 ]}\nobreak
\hyperdef{L}{X7F29407D834B102E}{}
{\noindent\textcolor{FuncColor}{$\triangleright$\enspace\texttt{RSGraphSpanningTree({\mdseries\slshape graph[, search{\textunderscore}type]})\index{RSGraphSpanningTree@\texttt{RSGraphSpanningTree}}
\label{RSGraphSpanningTree}
}\hfill{\scriptsize (operation)}}\\
\textbf{\indent Returns:\ }
A spanning tree of the graph. 



 This function calculates a spanning tree of the graph based on the edges. This
means it does not include and loops or cycles (see \texttt{RSGraphIsCycle} (\ref{RSGraphIsCycle}) for the definition of a cycle in an RSGraph). By default it uses a breadth
first search. 

 If the optional parameter \mbox{\texttt{\mdseries\slshape search{\textunderscore}type}} is equal to the string \texttt{"dfs"} then it will use a depth first search instead. If it is equal to the string \texttt{"bfs"} then it will use a breadth first search. If it is equal to any other string
then a warning will be raised and the function will use a breadth first
search. }

 
\begin{Verbatim}[commandchars=!@|,fontsize=\small,frame=single,label=Example]
  !gapprompt@gap>| !gapinput@adj_list := [[1, 2], [2, 1], [1, 3], [3, 1], [1, 4], [4, 1], \|
  !gapprompt@gap>| !gapinput@             [2, 3], [3, 2], [3, 4], [4, 3]];;|
  !gapprompt@gap>| !gapinput@rev_map := (1,2)(3,4)(5,6)(7,8)(9,10);;|
  !gapprompt@gap>| !gapinput@graph := RSGraphByAdjacencyList(adj_list, rev_map);;|
  !gapprompt@gap>| !gapinput@tree_bfs := RSGraphSpanningTree(graph);;|
  !gapprompt@gap>| !gapinput@Print(tree_bfs);|
  Vertices = { 1, 2, 3, 4 }
  Arcs = {
          1 = ( origin = 1, terminus = 2, inverse = 2 )
          2 = ( origin = 2, terminus = 1, inverse = 1 )
          3 = ( origin = 1, terminus = 3, inverse = 4 )
          4 = ( origin = 3, terminus = 1, inverse = 3 )
          5 = ( origin = 1, terminus = 4, inverse = 6 )
          6 = ( origin = 4, terminus = 1, inverse = 5 )
  }
  Reverse Map = (1,2)(3,4)(5,6)(7,8)(9,10)
  !gapprompt@gap>| !gapinput@tree_dfs := RSGraphSpanningTree(graph, "dfs");;|
  !gapprompt@gap>| !gapinput@Print(tree_dfs);|
  Vertices = { 1, 2, 3, 4 }
  Arcs = {
          1 = ( origin = 1, terminus = 2, inverse = 2 )
          2 = ( origin = 2, terminus = 1, inverse = 1 )
          7 = ( origin = 2, terminus = 3, inverse = 8 )
          8 = ( origin = 3, terminus = 2, inverse = 7 )
          9 = ( origin = 3, terminus = 4, inverse = 10 )
          10 = ( origin = 4, terminus = 3, inverse = 9 )
  }
  Reverse Map = (1,2)(3,4)(5,6)(7,8)(9,10)
\end{Verbatim}
 }

 }

   
\chapter{\textcolor{Chapter }{Local Action Diagrams}}\label{Chapter_Local_Action_Diagrams}
\logpage{[ 2, 0, 0 ]}
\hyperdef{L}{X7E4B6DDB802E1F7E}{}
{
  

 A local action diagram $\Delta = (\Gamma, (X_a), (G(v)))$ consists of: 
\begin{itemize}
\item  a graph $\Gamma$ (an \texttt{RSGraph}), 
\item  for each arc $a \in A(\Gamma)$a \emph{colour set} $|X_a|$ disjoint from every other colour set, 
\item  for each vertex $v \in V(\Gamma)$ a group $G(v)$ labelling the vertex such that $G(V)$ acts on $\bigsqcup_{a \in o^{-1}(v)}X_{a}$ and the orbits of $G(V)$ are precisely the colour sets $|X_a|$ for $a \in o^{-1}(v)$. 
\end{itemize}
 

 Local action diagrams are a powerful tool for studying $(P)$\texttt{\symbol{45}}closed groups because isomorphism classes of local action
diagrams are in a one\texttt{\symbol{45}}to\texttt{\symbol{45}}one
correspondence with conjugacy classes of $(P)$\texttt{\symbol{45}}closed groups. Furthermore, properties of the
corresponding $(P)$\texttt{\symbol{45}}closed group can be determined from the local action
diagram. In the case when a local action diagram is finite, we can represent
the diagram in \textsf{GAP} and determine properties of the (potentially infinite) corresponding group. 

 
\section{\textcolor{Chapter }{Creating Local Action Diagrams}}\label{Chapter_Local_Action_Diagrams_Section_Creating_Local_Action_Diagrams}
\logpage{[ 2, 1, 0 ]}
\hyperdef{L}{X84CBE9FE85CF5E9D}{}
{
  

 

\subsection{\textcolor{Chapter }{IsLocalActionDiagram (for an object)}}
\logpage{[ 2, 1, 1 ]}\nobreak
\hyperdef{L}{X79A1FE3686A8ACCD}{}
{\noindent\textcolor{FuncColor}{$\triangleright$\enspace\texttt{IsLocalActionDiagram({\mdseries\slshape obj})\index{IsLocalActionDiagram@\texttt{IsLocalActionDiagram}!for an object}
\label{IsLocalActionDiagram:for an object}
}\hfill{\scriptsize (Category)}}\\
\textbf{\indent Returns:\ }
\texttt{true} if the object is of the category \texttt{IsLocalActionDiagram} and false otherwise. 



 Every local action diagram belongs to the category \texttt{IsLocalActionDiagram}. Furthermore, every local action diagram is an immutable attribute storing
object. }

 

 

\subsection{\textcolor{Chapter }{LocalActionDiagramFromData (for an RSGraph, list or record of vertex labels, and list or record of arc labels)}}
\logpage{[ 2, 1, 2 ]}\nobreak
\label{LADFromData}
\hyperdef{L}{X7EB208B47BAB2BC1}{}
{\noindent\textcolor{FuncColor}{$\triangleright$\enspace\texttt{LocalActionDiagramFromData({\mdseries\slshape graph, vertex{\textunderscore}labels, arc{\textunderscore}labels})\index{LocalActionDiagramFromData@\texttt{LocalActionDiagramFromData}!for an RSGraph, list or record of vertex labels, and list or record of arc labels}
\label{LocalActionDiagramFromData:for an RSGraph, list or record of vertex labels, and list or record of arc labels}
}\hfill{\scriptsize (operation)}}\\
\noindent\textcolor{FuncColor}{$\triangleright$\enspace\texttt{LocalActionDiagramFromDataNC({\mdseries\slshape graph, vertex{\textunderscore}labels, arc{\textunderscore}labels})\index{LocalActionDiagramFromDataNC@\texttt{LocalActionDiagramFromDataNC}!for an RSGraph, list or record of vertex labels, and list or record of arc labels}
\label{LocalActionDiagramFromDataNC:for an RSGraph, list or record of vertex labels, and list or record of arc labels}
}\hfill{\scriptsize (operation)}}\\
\textbf{\indent Returns:\ }
A local action diagram object. 



 This function constructs a local action diagram by having all of the diagrams
data given to it. This is done by providing the \texttt{RSGraph} of the diagram and either lists of vertex and arc labels or records of vertex
and arc labels. 

 If two lists are provided then the vertex and arc labels are assumed to be in
the same order as \texttt{RSGraphVertices(\mbox{\texttt{\mdseries\slshape graph}})} and \texttt{RSGraphArcIDs(\mbox{\texttt{\mdseries\slshape graph}})} respectively. This means, for example, that element \texttt{idx} of \mbox{\texttt{\mdseries\slshape vertex{\textunderscore}labels}} is the label of the vertex with id \texttt{RSGraphVertexIDs(\mbox{\texttt{\mdseries\slshape graph}})[idx]}. If two records are provided then the names of the record elements are the
vertex and arc ids respectively and the components are their labels. 

 The vertex labels are permutation groups. The domain each label acts on is
stored in the mutable attribute \texttt{PermGroupDomain} (see \texttt{PermGroupDomain} (\ref{PermGroupDomain})). By default, this is equal to the \texttt{MovedPoints} of the group but can be changed to allow for the group to have fixed points. 

 The arc labels are lists of integers. Each arc must be labelled by by an orbit
of the vertex label at the vertex that arc originates at. Furthermore, every
vertex must have an arc originating at it for each orbit of its vertex label.
Note that the orbits of a vertex label \texttt{G} are determined by \texttt{Orbits(G, PermGroupDomain(G))}. 

 Unlike the formal definition of a local action diagram, we allow for overlap
between the domains of different vertex labels. For example, a two vertex
local action diagram could be labelled by $C_2$ and $S_3$ and they can have domains $\{1, 2\}$ and $\{1, 2, 3\}$ respectively. 

 The NC variant of the function does not check that there is a label for each
vertex and arc (and exactly this many labels), that the domains of the groups
match up with the arc labels, or that the orbits of the vertex labels match up
with the arc labels. }

 
\begin{Verbatim}[commandchars=!@|,fontsize=\small,frame=single,label=Example]
  !gapprompt@gap>| !gapinput@graph := RSGraphByAdjacencyList([[1, 1], [1, 2], [2, 1]], (2,3));;|
  !gapprompt@gap>| !gapinput@label_1 := Group((1,2));;|
  !gapprompt@gap>| !gapinput@SetPermGroupDomain(label_1, [1, 2, 3]);;|
  !gapprompt@gap>| !gapinput@label_2 := Group((1,2,3));;|
  !gapprompt@gap>| !gapinput@vertex_labels := [label_1, label_2];|
  !gapprompt@gap>| !gapinput@arc_labels := [[1, 2], [3], [1,2,3]];|
  !gapprompt@gap>| !gapinput@lad := LocalActionDiagramFromData(graph, vertex_labels, arc_labels);|
  <LocalActionDiagram with 2 vertices and 3 arcs>
  !gapprompt@gap>| !gapinput@Print(lad);|
  Vertices = { 1, 2 }
  Arcs = {
          1 = ( origin = 1, terminus = 1, inverse = 1 )
          2 = ( origin = 1, terminus = 2, inverse = 3 )
          3 = ( origin = 2, terminus = 1, inverse = 2 )
  }
  Reverse Map = (2,3)
  Vertex Labels = {
          1 = Group( [ (1,2) ] )
          2 = Group( [ (1,2,3) ] )
  }
  Arc Labels = {
          1 = [ 1, 2 ]
          2 = [ 3 ]
          3 = [ 1, 2, 3 ]
  }
\end{Verbatim}
 

\subsection{\textcolor{Chapter }{LocalActionDiagramFromBurgerMozesUniversalGroup}}
\logpage{[ 2, 1, 3 ]}\nobreak
\hyperdef{L}{X785DF9A07E409955}{}
{\noindent\textcolor{FuncColor}{$\triangleright$\enspace\texttt{LocalActionDiagramFromBurgerMozesUniversalGroup({\mdseries\slshape perm{\textunderscore}group})\index{LocalActionDiagramFromBurgerMozesUniversalGroup@\texttt{Local}\-\texttt{Action}\-\texttt{Diagram}\-\texttt{From}\-\texttt{Burger}\-\texttt{Mozes}\-\texttt{Universal}\-\texttt{Group}}
\label{LocalActionDiagramFromBurgerMozesUniversalGroup}
}\hfill{\scriptsize (operation)}}\\
\textbf{\indent Returns:\ }
A local action diagram object. 



 Given a permutation group $F$ (\mbox{\texttt{\mdseries\slshape perm{\textunderscore}group}}) this function returns the local action diagram corresponding to the
Burger\texttt{\symbol{45}}Mozes universal group $U(F)$ (see [cite Burger\texttt{\symbol{45}}Mozes and Colin and Simon for
construction?]). This is a local action diagram with one vertex labelled by $F$ and a self\texttt{\symbol{45}}reverse loop for each orbit of $F$. Note that $F$ acts on the domain \texttt{PermGroupDomain(\mbox{\texttt{\mdseries\slshape perm{\textunderscore}group}})}. }

 
\begin{Verbatim}[commandchars=!@|,fontsize=\small,frame=single,label=Example]
  !gapprompt@gap>| !gapinput@lad := LocalActionDiagramFromBurgerMozesUniversalGroup(Group((1,2)(3,4)));|
  <U(Group( [ (1,2)(3,4) ] )) (as a Local Action Diagram)>
\end{Verbatim}
 

\subsection{\textcolor{Chapter }{LocalActionDiagramConstructBipartitionPreserving}}
\logpage{[ 2, 1, 4 ]}\nobreak
\hyperdef{L}{X81C137048009714A}{}
{\noindent\textcolor{FuncColor}{$\triangleright$\enspace\texttt{LocalActionDiagramConstructBipartitionPreserving({\mdseries\slshape lad})\index{LocalActionDiagramConstructBipartitionPreserving@\texttt{Local}\-\texttt{Action}\-\texttt{Diagram}\-\texttt{Construct}\-\texttt{Bipartition}\-\texttt{Preserving}}
\label{LocalActionDiagramConstructBipartitionPreserving}
}\hfill{\scriptsize (operation)}}\\
\textbf{\indent Returns:\ }
A local action diagram object. 



 Given a single vertex local action diagram $\Delta$ (\mbox{\texttt{\mdseries\slshape lad}}) corresponding to the group $U(\Delta)$ acting on a tree $T$ this function returns the local action diagram corresponding to the group $U(\Delta)^{\circ} = \{g \in U(\Delta)\, |\, \forall\, v \in V(T),\, d(v, gv)
\equiv 0 \pmod{2}\}$. Intuitively this is the group of all automorphisms the preserve the
bipartition of $T$ which always exists since a single vertex local action diagram is vertex
transitive. 

 Given a group $G$ labelling the single vertex of $\Delta$ and arc labels $X_a$ $r$ this local action diagram is constructed as follows: 
\begin{itemize}
\item  Make a graph with two vertices and an edge for each loop of the graph for $\Delta$. 
\item  Label both vertices by $G$. 
\item  For each arc $a$ in the graph of $\Delta$ label the corresponding arc in the new graph by $X_{a}$ and the reverse of this arc by $X_{\overline{a}}$. 
\end{itemize}
 }

 }

 
\section{\textcolor{Chapter }{Local Action Diagram Attributes and Properties}}\label{Chapter_Local_Action_Diagrams_Section_Local_Action_Diagram_Attributes_and_Properties}
\logpage{[ 2, 2, 0 ]}
\hyperdef{L}{X8048984C81A6C0C2}{}
{
  

\subsection{\textcolor{Chapter }{LocalActionDiagramRSGraph}}
\logpage{[ 2, 2, 1 ]}\nobreak
\hyperdef{L}{X7D2ED3BD7CF95FA4}{}
{\noindent\textcolor{FuncColor}{$\triangleright$\enspace\texttt{LocalActionDiagramRSGraph({\mdseries\slshape lad})\index{LocalActionDiagramRSGraph@\texttt{LocalActionDiagramRSGraph}}
\label{LocalActionDiagramRSGraph}
}\hfill{\scriptsize (attribute)}}\\
\textbf{\indent Returns:\ }
RSGraph associated with the local action diagram. 



 

 }

 

\subsection{\textcolor{Chapter }{LocalActionDiagramVertexLabels}}
\logpage{[ 2, 2, 2 ]}\nobreak
\hyperdef{L}{X7D9D2615812AE56B}{}
{\noindent\textcolor{FuncColor}{$\triangleright$\enspace\texttt{LocalActionDiagramVertexLabels({\mdseries\slshape lad})\index{LocalActionDiagramVertexLabels@\texttt{LocalActionDiagramVertexLabels}}
\label{LocalActionDiagramVertexLabels}
}\hfill{\scriptsize (attribute)}}\\
\textbf{\indent Returns:\ }
Vertex labels of the local action diagram. 



 

 }

 

\subsection{\textcolor{Chapter }{LocalActionDiagramArcLabels}}
\logpage{[ 2, 2, 3 ]}\nobreak
\hyperdef{L}{X7FA7E80680DE84C2}{}
{\noindent\textcolor{FuncColor}{$\triangleright$\enspace\texttt{LocalActionDiagramArcLabels({\mdseries\slshape lad})\index{LocalActionDiagramArcLabels@\texttt{LocalActionDiagramArcLabels}}
\label{LocalActionDiagramArcLabels}
}\hfill{\scriptsize (attribute)}}\\
\textbf{\indent Returns:\ }
Arc labels of the local action diagram. 



 

 }

 

\subsection{\textcolor{Chapter }{LocalActionDiagramVertices}}
\logpage{[ 2, 2, 4 ]}\nobreak
\hyperdef{L}{X874DED9479AE6960}{}
{\noindent\textcolor{FuncColor}{$\triangleright$\enspace\texttt{LocalActionDiagramVertices({\mdseries\slshape lad})\index{LocalActionDiagramVertices@\texttt{LocalActionDiagramVertices}}
\label{LocalActionDiagramVertices}
}\hfill{\scriptsize (attribute)}}\\
\textbf{\indent Returns:\ }
The vertices of the local action diagrams graph. 



 Shorthand for \texttt{RSGraphVertices(LocalActionDiagramRSGraph(\mbox{\texttt{\mdseries\slshape lad}}))}. }

 

\subsection{\textcolor{Chapter }{LocalActionDiagramArcs}}
\logpage{[ 2, 2, 5 ]}\nobreak
\hyperdef{L}{X7F304E9E7D9876A7}{}
{\noindent\textcolor{FuncColor}{$\triangleright$\enspace\texttt{LocalActionDiagramArcs({\mdseries\slshape lad})\index{LocalActionDiagramArcs@\texttt{LocalActionDiagramArcs}}
\label{LocalActionDiagramArcs}
}\hfill{\scriptsize (attribute)}}\\
\textbf{\indent Returns:\ }
The arcs of the local action diagrams graph. 



 Shorthand for \texttt{RSGraphArcs(LocalActionDiagramRSGraph(\mbox{\texttt{\mdseries\slshape lad}}))}. }

 

\subsection{\textcolor{Chapter }{LocalActionDiagramArcIDs}}
\logpage{[ 2, 2, 6 ]}\nobreak
\hyperdef{L}{X823F306F799E4F11}{}
{\noindent\textcolor{FuncColor}{$\triangleright$\enspace\texttt{LocalActionDiagramArcIDs({\mdseries\slshape lad})\index{LocalActionDiagramArcIDs@\texttt{LocalActionDiagramArcIDs}}
\label{LocalActionDiagramArcIDs}
}\hfill{\scriptsize (attribute)}}\\
\textbf{\indent Returns:\ }
The arc ids of the local action diagrams graph. 



 Shorthand for \texttt{RSGraphArcIDs(LocalActionDiagramRSGraph(\mbox{\texttt{\mdseries\slshape lad}}))}. }

 

\subsection{\textcolor{Chapter }{LocalActionDiagramNumberVertices}}
\logpage{[ 2, 2, 7 ]}\nobreak
\hyperdef{L}{X7F687E2E879DE15C}{}
{\noindent\textcolor{FuncColor}{$\triangleright$\enspace\texttt{LocalActionDiagramNumberVertices({\mdseries\slshape lad})\index{LocalActionDiagramNumberVertices@\texttt{LocalActionDiagramNumberVertices}}
\label{LocalActionDiagramNumberVertices}
}\hfill{\scriptsize (attribute)}}\\
\textbf{\indent Returns:\ }
The number of vertices in the local action diagrams graph. 



 Shorthand for \texttt{RSGraphNumberVertices(LocalActionDiagramRSGraph(\mbox{\texttt{\mdseries\slshape lad}}))}. }

 

\subsection{\textcolor{Chapter }{LocalActionDiagramNumberArcs}}
\logpage{[ 2, 2, 8 ]}\nobreak
\hyperdef{L}{X792FCDFB843B6A7A}{}
{\noindent\textcolor{FuncColor}{$\triangleright$\enspace\texttt{LocalActionDiagramNumberArcs({\mdseries\slshape lad})\index{LocalActionDiagramNumberArcs@\texttt{LocalActionDiagramNumberArcs}}
\label{LocalActionDiagramNumberArcs}
}\hfill{\scriptsize (attribute)}}\\
\textbf{\indent Returns:\ }
The number of arcs in the local action diagrams graph. 



 Shorthand for \texttt{RSGraphNumberArcs(LocalActionDiagramRSGraph(\mbox{\texttt{\mdseries\slshape lad}}))}. }

 

\subsection{\textcolor{Chapter }{LocalActionDiagramReverseMap}}
\logpage{[ 2, 2, 9 ]}\nobreak
\hyperdef{L}{X7F74CB4E7D5078DA}{}
{\noindent\textcolor{FuncColor}{$\triangleright$\enspace\texttt{LocalActionDiagramReverseMap({\mdseries\slshape lad})\index{LocalActionDiagramReverseMap@\texttt{LocalActionDiagramReverseMap}}
\label{LocalActionDiagramReverseMap}
}\hfill{\scriptsize (attribute)}}\\
\textbf{\indent Returns:\ }
The reverse map of the local action diagrams graph. 



 Shorthand for \texttt{RSGraphReverseMap(LocalActionDiagramRSGraph(\mbox{\texttt{\mdseries\slshape lad}}))}. }

 

\subsection{\textcolor{Chapter }{LocalActionDiagramOutNeighbours}}
\logpage{[ 2, 2, 10 ]}\nobreak
\hyperdef{L}{X825FA59783951DCE}{}
{\noindent\textcolor{FuncColor}{$\triangleright$\enspace\texttt{LocalActionDiagramOutNeighbours({\mdseries\slshape lad})\index{LocalActionDiagramOutNeighbours@\texttt{LocalActionDiagramOutNeighbours}}
\label{LocalActionDiagramOutNeighbours}
}\hfill{\scriptsize (attribute)}}\\
\textbf{\indent Returns:\ }
The record of out neighbours of the local action diagrams graph. 



 Shorthand for \texttt{RSGraphOutNeighbours(LocalActionDiagramRSGraph(\mbox{\texttt{\mdseries\slshape lad}}))}. }

 

\subsection{\textcolor{Chapter }{LocalActionDiagramInNeighbours}}
\logpage{[ 2, 2, 11 ]}\nobreak
\hyperdef{L}{X82253261848258F8}{}
{\noindent\textcolor{FuncColor}{$\triangleright$\enspace\texttt{LocalActionDiagramInNeighbours({\mdseries\slshape lad})\index{LocalActionDiagramInNeighbours@\texttt{LocalActionDiagramInNeighbours}}
\label{LocalActionDiagramInNeighbours}
}\hfill{\scriptsize (attribute)}}\\
\textbf{\indent Returns:\ }
The record of in neighbours of the local action diagrams graph. 



 Shorthand for \texttt{RSGraphInNeighbours(LocalActionDiagramRSGraph(\mbox{\texttt{\mdseries\slshape lad}}))}. }

 

\subsection{\textcolor{Chapter }{LocalActionDiagramOutArcs}}
\logpage{[ 2, 2, 12 ]}\nobreak
\hyperdef{L}{X792BBA5B7DEB9D04}{}
{\noindent\textcolor{FuncColor}{$\triangleright$\enspace\texttt{LocalActionDiagramOutArcs({\mdseries\slshape lad})\index{LocalActionDiagramOutArcs@\texttt{LocalActionDiagramOutArcs}}
\label{LocalActionDiagramOutArcs}
}\hfill{\scriptsize (attribute)}}\\
\textbf{\indent Returns:\ }
The record of out arcs of the local action diagrams graph. 



 Shorthand for \texttt{RSGraphOutArcs(LocalActionDiagramRSGraph(\mbox{\texttt{\mdseries\slshape lad}}))}. }

 

\subsection{\textcolor{Chapter }{LocalActionDiagramInArcs}}
\logpage{[ 2, 2, 13 ]}\nobreak
\hyperdef{L}{X79F8BECC81804228}{}
{\noindent\textcolor{FuncColor}{$\triangleright$\enspace\texttt{LocalActionDiagramInArcs({\mdseries\slshape lad})\index{LocalActionDiagramInArcs@\texttt{LocalActionDiagramInArcs}}
\label{LocalActionDiagramInArcs}
}\hfill{\scriptsize (attribute)}}\\
\textbf{\indent Returns:\ }
The record of in arcs of the local action diagrams graph. 



 Shorthand for \texttt{RSGraphInArcs(LocalActionDiagramRSGraph(\mbox{\texttt{\mdseries\slshape lad}}))}. }

 

\subsection{\textcolor{Chapter }{LocalActionDiagramGroupName}}
\logpage{[ 2, 2, 14 ]}\nobreak
\hyperdef{L}{X7978DED887F2496F}{}
{\noindent\textcolor{FuncColor}{$\triangleright$\enspace\texttt{LocalActionDiagramGroupName({\mdseries\slshape lad})\index{LocalActionDiagramGroupName@\texttt{LocalActionDiagramGroupName}}
\label{LocalActionDiagramGroupName}
}\hfill{\scriptsize (attribute)}}\\
\textbf{\indent Returns:\ }
A name for the associated universal group of the local action diagram. 



 This is a mutable attribute which stores a human readable name for the
associated universal group. By default it is equal to the empty string. In
this case the \texttt{View} function will print \texttt{{\textless}LocalActionDiagram with n vertices and m arcs{\textgreater}} where \texttt{n} is the number of vertices and \texttt{m} is the number of arcs. If this attribute is set to to any other string then
the \texttt{View} function will print \texttt{{\textless}\texttt{\symbol{123}}group{\textunderscore}name\texttt{\symbol{125}}
(as a Local Action Diagram){\textgreater}}. }

 
\begin{Verbatim}[commandchars=!@|,fontsize=\small,frame=single,label=Example]
  !gapprompt@gap>| !gapinput@graph := RSGraphByAdjacencyList([[1, 1]], ());;|
  !gapprompt@gap>| !gapinput@vertex_labels := [SymmetricGroup(3)];;|
  !gapprompt@gap>| !gapinput@arc_labels := [[1,2,3]];;|
  !gapprompt@gap>| !gapinput@lad := LocalActionDiagramFromData(graph, vertex_labels, arc_labels);|
  <LocalActionDiagram with 1 vertex and 1 arc>
  !gapprompt@gap>| !gapinput@SetLocalActionDiagramGroupName(lad, "Aut(T_3)");|
  !gapprompt@gap>| !gapinput@View(lad);|
  <Aut(T_3) (as a Local Action Diagram)>
\end{Verbatim}
 

\subsection{\textcolor{Chapter }{LocalActionDiagramRegularTree}}
\logpage{[ 2, 2, 15 ]}\nobreak
\hyperdef{L}{X7F2B374A8049A165}{}
{\noindent\textcolor{FuncColor}{$\triangleright$\enspace\texttt{LocalActionDiagramRegularTree({\mdseries\slshape lad})\index{LocalActionDiagramRegularTree@\texttt{LocalActionDiagramRegularTree}}
\label{LocalActionDiagramRegularTree}
}\hfill{\scriptsize (attribute)}}\\
\textbf{\indent Returns:\ }
The degree of the regular tree the universal group acts on or \texttt{fail} if it does not act on one. 



 If every vertex label of the local action diagram has a domain of size $d$ then the associated universal group acts on a regular tree of degree $d$. This functions returns $d$ if this is the case and \texttt{fail} otherwise. }

 
\begin{Verbatim}[commandchars=!@|,fontsize=\small,frame=single,label=Example]
  !gapprompt@gap>| !gapinput@graph := RSGraphByAdjacencyList([[1, 2], [2, 1]], (1,2));;|
  !gapprompt@gap>| !gapinput@vertex_labels := [SymmetricGroup(3), Group((1,2,3))];;|
  !gapprompt@gap>| !gapinput@arc_labels := [[1,2,3], [1,2,3]];;|
  !gapprompt@gap>| !gapinput@lad := LocalActionDiagramFromData(graph, vertex_labels, arc_labels);;|
  !gapprompt@gap>| !gapinput@LocalActionDiagramRegularTree(lad);|
  3
  !gapprompt@gap>| !gapinput@vertex_labels_2 := [SymmetricGroup(3), Group((1,2))];;|
  !gapprompt@gap>| !gapinput@arc_labels_2 := [[1,2,3], [1,2]];;|
  !gapprompt@gap>| !gapinput@lad_2 := LocalActionDiagramFromData(graph, vertex_labels_2, arc_labels_2);;|
  !gapprompt@gap>| !gapinput@LocalActionDiagramRegularTree(lad_2);|
  fail
\end{Verbatim}
 

\subsection{\textcolor{Chapter }{LocalActionDiagramScopos}}
\logpage{[ 2, 2, 16 ]}\nobreak
\hyperdef{L}{X78AB8E5D83D7823D}{}
{\noindent\textcolor{FuncColor}{$\triangleright$\enspace\texttt{LocalActionDiagramScopos({\mdseries\slshape lad})\index{LocalActionDiagramScopos@\texttt{LocalActionDiagramScopos}}
\label{LocalActionDiagramScopos}
}\hfill{\scriptsize (attribute)}}\\
\textbf{\indent Returns:\ }
The list of scopos of the local action diagram. 



 (CITE REID SMITH HERE) Let $\Delta = (\Gamma, (X_a), (G(v)))$ be a local action diagram. A strongly confluent partial orientation (\emph{scopo}) of $\Delta$ is a subset $O$ of $A(\Gamma)$ such that: 
\begin{itemize}
\item  if $a \in O$ then $\overline{a} \notin O$, 
\item  if $a \in O$ then $\left|X_{a}\right| = 1$, and 
\item  for all $v \in V(\Gamma)$ if $O$ contains an arc $a$ originating at $v$ then $O$ contains all arcs that terminate at $v$ except for $\overline{a}$. 
\end{itemize}
 

 This function performs an iterative search for all scopos of a local action
diagram. It stores each scopo as a list of arc ids in the scopo. Note that the
empty set is a scopo and this is represented as the empty list. }

 
\begin{Verbatim}[commandchars=!@|,fontsize=\small,frame=single,label=Example]
  !gapprompt@gap>| !gapinput@graph := RSGraphByAdjacencyList ([[1, 2], [2, 1]], (1,2));;|
  !gapprompt@gap>| !gapinput@group := Group(());;|
  !gapprompt@gap>| !gapinput@SetPermGroupDomain(group, [1]);;|
  !gapprompt@gap>| !gapinput@vertex_labels := [group, group];;|
  !gapprompt@gap>| !gapinput@arc_labels := [[1], [1]];;|
  !gapprompt@gap>| !gapinput@lad := LocalActionDiagramFromData(graph, vertex_labels, arc_labels);;|
  !gapprompt@gap>| !gapinput@LocalActionDiagramScopos(lad);|
  [ [ ], [ 1 ], [ 2 ]]
\end{Verbatim}
 

\subsection{\textcolor{Chapter }{LocalActionDiagramGroupType}}
\logpage{[ 2, 2, 17 ]}\nobreak
\hyperdef{L}{X84570C4A84D97EB2}{}
{\noindent\textcolor{FuncColor}{$\triangleright$\enspace\texttt{LocalActionDiagramGroupType({\mdseries\slshape lad})\index{LocalActionDiagramGroupType@\texttt{LocalActionDiagramGroupType}}
\label{LocalActionDiagramGroupType}
}\hfill{\scriptsize (attribute)}}\\
\textbf{\indent Returns:\ }
The type of the corresponding universal group. 



 (CITE REID SMITH HERE) Every group acting on a tree can be split into one of
six mutually exclusive types: \emph{fixed vertex}, \emph{edge inversion}, \emph{lineal}, \emph{horocyclic}, \emph{focal}, and \emph{general}. This types can be recognised from the groups corresponding local action
diagram by analysing the scopos of the local action diagram (see CITE REID
SMITH). 

 This function first calculates the scopos of the local action diagram (see \texttt{LocalActionDiagramScopos} (\ref{LocalActionDiagramScopos})) and then uses this to determine the type of the corresponding group. It
outputs this types as a string. }

 
\begin{Verbatim}[commandchars=!@|,fontsize=\small,frame=single,label=Example]
  !gapprompt@gap>| !gapinput@graph := RSGraphByAdjacencyList ([[1, 2], [2, 1]], (1,2));;|
  !gapprompt@gap>| !gapinput@group := Group(());;|
  !gapprompt@gap>| !gapinput@SetPermGroupDomain(group, [1]);;|
  !gapprompt@gap>| !gapinput@vertex_labels := [group, group];;|
  !gapprompt@gap>| !gapinput@arc_labels := [[1], [1]];;|
  !gapprompt@gap>| !gapinput@lad := LocalActionDiagramFromData(graph, vertex_labels, arc_labels);;|
  !gapprompt@gap>| !gapinput@LocalActionDiagramGroupType(lad);|
  "General"
\end{Verbatim}
 

\subsection{\textcolor{Chapter }{LocalActionDiagramIsDiscrete}}
\logpage{[ 2, 2, 18 ]}\nobreak
\hyperdef{L}{X83CEFD1682867DD9}{}
{\noindent\textcolor{FuncColor}{$\triangleright$\enspace\texttt{LocalActionDiagramIsDiscrete({\mdseries\slshape lad})\index{LocalActionDiagramIsDiscrete@\texttt{LocalActionDiagramIsDiscrete}}
\label{LocalActionDiagramIsDiscrete}
}\hfill{\scriptsize (property)}}\\
\textbf{\indent Returns:\ }
\texttt{true} if the corresponding universal group is discrete and false otherwise. 



 This function implements (CITE THEOREM HERE) to determine if the corresponding
universal group is discrete. This requires calculation of the diagrams group
type (see \texttt{LocalActionDiagramGroupType} (\ref{LocalActionDiagramGroupType})). Furthermore, if the group is discrete then it is also uniscalar and
unimodular and so if this function returns \texttt{true} then it also sets \texttt{LocalActionDiagramIsUniscalar} (\ref{LocalActionDiagramIsUniscalar}) and \texttt{LocalActionDiagramIsUnimodular} (\ref{LocalActionDiagramIsUnimodular}) to true. }

 
\begin{Verbatim}[commandchars=!@|,fontsize=\small,frame=single,label=Example]
  !gapprompt@gap>| !gapinput@graph := RSGraphByAdjacencyList ([[1, 2], [2, 1]], (1,2));;|
  !gapprompt@gap>| !gapinput@group := Group(());;|
  !gapprompt@gap>| !gapinput@SetPermGroupDomain(group, [1]);;|
  !gapprompt@gap>| !gapinput@vertex_labels := [group, group];;|
  !gapprompt@gap>| !gapinput@arc_labels := [[1], [1]];;|
  !gapprompt@gap>| !gapinput@lad := LocalActionDiagramFromData(graph, vertex_labels, arc_labels);;|
  !gapprompt@gap>| !gapinput@LocalActionDiagramIsDiscrete(lad);|
  true
\end{Verbatim}
 

\subsection{\textcolor{Chapter }{LocalActionDiagramIsUniscalar}}
\logpage{[ 2, 2, 19 ]}\nobreak
\hyperdef{L}{X82155DA17AAC826E}{}
{\noindent\textcolor{FuncColor}{$\triangleright$\enspace\texttt{LocalActionDiagramIsUniscalar({\mdseries\slshape lad})\index{LocalActionDiagramIsUniscalar@\texttt{LocalActionDiagramIsUniscalar}}
\label{LocalActionDiagramIsUniscalar}
}\hfill{\scriptsize (property)}}\\
\textbf{\indent Returns:\ }
\texttt{true} if the corresponding universal group is uniscalar and false otherwise. 



 This function implements (CITE THEOREM HERE) to determine if the corresponding
universal group is uniscalar. This requires calculation of the diagrams group
type (see \texttt{LocalActionDiagramGroupType} (\ref{LocalActionDiagramGroupType})). Furthermore, if the group is uniscalar then it is also unimodular so if
this function returns \texttt{true} then it also sets \texttt{LocalActionDiagramIsUnimodular} (\ref{LocalActionDiagramIsUnimodular}) to true. }

 
\begin{Verbatim}[commandchars=!@|,fontsize=\small,frame=single,label=Example]
  !gapprompt@gap>| !gapinput@graph := RSGraphByAdjacencyList ([[1, 2], [2, 1], [2, 2]], (1,2));;|
  !gapprompt@gap>| !gapinput@group := Group(());;|
  !gapprompt@gap>| !gapinput@group_2 := Group((1, 2), (1,2,3), (4,5));;|
  !gapprompt@gap>| !gapinput@SetPermGroupDomain(group, [1]);;|
  !gapprompt@gap>| !gapinput@vertex_labels := [group, group_2];;|
  !gapprompt@gap>| !gapinput@arc_labels := [[1], [1, 2, 3], [4, 5]];;|
  !gapprompt@gap>| !gapinput@lad := LocalActionDiagramFromData(graph, vertex_labels, arc_labels);;|
  !gapprompt@gap>| !gapinput@LocalActionDiagramIsDiscrete(lad);|
  false
  !gapprompt@gap>| !gapinput@LocalActionDiagramIsUniscalar(lad);|
  true
\end{Verbatim}
 

\subsection{\textcolor{Chapter }{LocalActionDiagramIsUnimodular}}
\logpage{[ 2, 2, 20 ]}\nobreak
\hyperdef{L}{X816A48F4838326B8}{}
{\noindent\textcolor{FuncColor}{$\triangleright$\enspace\texttt{LocalActionDiagramIsUnimodular({\mdseries\slshape lad})\index{LocalActionDiagramIsUnimodular@\texttt{LocalActionDiagramIsUnimodular}}
\label{LocalActionDiagramIsUnimodular}
}\hfill{\scriptsize (property)}}\\
\textbf{\indent Returns:\ }
\texttt{true} if the corresponding universal group is unimodular and false otherwise. 



 This function implements (CITE THEOREM HERE both us and BK) to determine if
the corresponding universal group is unimodular. This requires calculation of
a spanning tree of the graph (see \texttt{RSGraphSpanningTree} (\ref{RSGraphSpanningTree})). }

 
\begin{Verbatim}[commandchars=!@|,fontsize=\small,frame=single,label=Example]
  !gapprompt@gap>| !gapinput@adj_list := [[1, 2], [2, 1], [2, 3], [3, 2], [3, 1], [1, 3], [1, 1]];;|
  !gapprompt@gap>| !gapinput@graph := RSGraphByAdjacencyList(adj_list, (1,2)(3,4)(5,6));;|
  !gapprompt@gap>| !gapinput@vertex_labels := [];;|
  !gapprompt@gap>| !gapinput@Add(vertex_labels, Group((1, 2, 3), (4,5), (6,7)));;|
  !gapprompt@gap>| !gapinput@Add(vertex_labels, Group((1, 2), (3, 4, 5)));;|
  !gapprompt@gap>| !gapinput@Add(vertex_labels, Group((1, 2), (3, 4)));;|
  !gapprompt@gap>| !gapinput@arc_labels := [[1, 2, 3], [3, 4, 5], [1, 2], [1, 2], [3,4], [4,5], [6,7]];;|
  !gapprompt@gap>| !gapinput@lad := LocalActionDiagramFromData(graph, vertex_labels, arc_labels);;|
  !gapprompt@gap>| !gapinput@LocalActionDiagramIsDiscrete(lad);|
  false
  !gapprompt@gap>| !gapinput@LocalActionDiagramIsUniscalar(lad);|
  false
  !gapprompt@gap>| !gapinput@LocalActionDiagramIsUnimodular(lad);|
  true
\end{Verbatim}
 

\subsection{\textcolor{Chapter }{LocalActionDiagramIsBurgerMozesUniversalGroup}}
\logpage{[ 2, 2, 21 ]}\nobreak
\hyperdef{L}{X87E1AE8587DA1FB1}{}
{\noindent\textcolor{FuncColor}{$\triangleright$\enspace\texttt{LocalActionDiagramIsBurgerMozesUniversalGroup({\mdseries\slshape lad})\index{LocalActionDiagramIsBurgerMozesUniversalGroup@\texttt{Local}\-\texttt{Action}\-\texttt{Diagram}\-\texttt{Is}\-\texttt{Burger}\-\texttt{Mozes}\-\texttt{Universal}\-\texttt{Group}}
\label{LocalActionDiagramIsBurgerMozesUniversalGroup}
}\hfill{\scriptsize (property)}}\\
\textbf{\indent Returns:\ }
\texttt{true} if the corresponding universal group is a Burger\texttt{\symbol{45}}Mozes
Universal Group and \texttt{false} otherwise. 



 If a local action diagram has a single vertex and all loops at the vertex are
self\texttt{\symbol{45}}reverse then the diagram corresponds to a
Burger\texttt{\symbol{45}}Mozes group (see PROBABLY SECOND COLIN SIMON PAPER).
If the given diagram has such a structure then this function returns true. It
also sets \texttt{LocalActionDiagramGroupName} (\ref{LocalActionDiagramGroupName}) to \texttt{"U(\texttt{\symbol{123}}vertex{\textunderscore}label\texttt{\symbol{125}})"} where \texttt{\texttt{\symbol{123}}vertex{\textunderscore}label\texttt{\symbol{125}}} is the name of the group labelling the vertex. }

 
\begin{Verbatim}[commandchars=!@|,fontsize=\small,frame=single,label=Example]
  !gapprompt@gap>| !gapinput@graph := RSGraphByAdjacencyList([[1, 1], [1, 1]], ());;|
  !gapprompt@gap>| !gapinput@vertex_labels := [Group((1, 2), (3, 4))];;|
  !gapprompt@gap>| !gapinput@SetName(vertex_labels[1], "C2 x C2");|
  !gapprompt@gap>| !gapinput@arc_labels := [[1, 2], [3, 4]];;|
  !gapprompt@gap>| !gapinput@lad := LocalActionDiagramFromData(graph, vertex_labels, arc_labels);;|
  !gapprompt@gap>| !gapinput@View(lad);|
  <LocalActionDiagram with 1 vertex and 2 arcs>
  !gapprompt@gap>| !gapinput@LocalActionDiagramIsBurgerMozesUniversalGroup(lad);|
  true
  !gapprompt@gap>| !gapinput@View(lad);|
  <U(C2 x C2) (as a Local Action Diagram)>
\end{Verbatim}
 

\subsection{\textcolor{Chapter }{LocalActionDiagramIsEndStabiliser}}
\logpage{[ 2, 2, 22 ]}\nobreak
\hyperdef{L}{X8443CE477DC98586}{}
{\noindent\textcolor{FuncColor}{$\triangleright$\enspace\texttt{LocalActionDiagramIsEndStabiliser({\mdseries\slshape lad})\index{LocalActionDiagramIsEndStabiliser@\texttt{LocalActionDiagramIsEndStabiliser}}
\label{LocalActionDiagramIsEndStabiliser}
}\hfill{\scriptsize (property)}}\\
\textbf{\indent Returns:\ }
\texttt{true} if the corresponding universal group is the stabiliser of a single end in the
full automorphism group of a regular tree and \texttt{false} otherwise. otherwise. 



 If a local action diagram has a single vertex with two mutually reverse loops
and the vertex is labelled by $S_{d-1}$ acting on $\{1, 2, \dots, d\}$ then the universal group is isomorphic to $\operatorname{Aut}(T_{d})_{\omega}$ where $\omega$ is an end of $T_{d}$. This function returns \texttt{true} if the group satisfies this structure and \texttt{false} otherwise. It also sets \texttt{LocalActionDiagramGroupName} (\ref{LocalActionDiagramGroupName}) to \texttt{Aut(Td){\textunderscore}omega}. }

 
\begin{Verbatim}[commandchars=!@|,fontsize=\small,frame=single,label=Example]
  !gapprompt@gap>| !gapinput@graph := RSGraphByAdjacencyList([[1, 1], [1, 1]], (1, 2));;|
  !gapprompt@gap>| !gapinput@vertex_labels := [SymmetricGroup(2)];;|
  !gapprompt@gap>| !gapinput@SetPermGroupDomain(vertex_labels[1], [1,2,3]);;|
  !gapprompt@gap>| !gapinput@arc_labels := [[1, 2], [3]];;|
  !gapprompt@gap>| !gapinput@lad := LocalActionDiagramFromData(graph, vertex_labels, arc_labels);;|
  !gapprompt@gap>| !gapinput@View(lad);|
  <LocalActionDiagram with 1 vertex and 2 arcs>
  !gapprompt@gap>| !gapinput@LocalActionDiagramIsEndStabiliser(lad);|
  true
  !gapprompt@gap>| !gapinput@View(lad);|
  <Aut(T3)_omega (as a Local Action Diagram)>
\end{Verbatim}
 

\subsection{\textcolor{Chapter }{LocalActionDiagramIsRegularTreeAutomorphismGroup}}
\logpage{[ 2, 2, 23 ]}\nobreak
\hyperdef{L}{X829A07D678914E89}{}
{\noindent\textcolor{FuncColor}{$\triangleright$\enspace\texttt{LocalActionDiagramIsRegularTreeAutomorphismGroup({\mdseries\slshape lad})\index{LocalActionDiagramIsRegularTreeAutomorphismGroup@\texttt{Local}\-\texttt{Action}\-\texttt{Diagram}\-\texttt{Is}\-\texttt{Regular}\-\texttt{Tree}\-\texttt{Automorphism}\-\texttt{Group}}
\label{LocalActionDiagramIsRegularTreeAutomorphismGroup}
}\hfill{\scriptsize (property)}}\\
\textbf{\indent Returns:\ }
\texttt{true} if the corresponding universal group is the full automorphism group of a
regular tree and \texttt{false} otherwise. 



 If a local action diagram has a single vertex with one
self\texttt{\symbol{45}}reverse loop and the vertex is labelled by $S_d$ then the universal group is isomorphic to $\operatorname{Aut}(T_d)$. This function returns \texttt{true} if the group satisfies this structure and \texttt{false} otherwise. It also sets \texttt{LocalActionDiagramGroupName} (\ref{LocalActionDiagramGroupName}) to \texttt{Aut(Td)}. }

 
\begin{Verbatim}[commandchars=!@|,fontsize=\small,frame=single,label=Example]
  !gapprompt@gap>| !gapinput@graph := RSGraphByAdjacencyList([[1, 1]], ());;|
  !gapprompt@gap>| !gapinput@vertex_labels := [SymmetricGroup(5)];;|
  !gapprompt@gap>| !gapinput@arc_labels := [[1..5]];;|
  !gapprompt@gap>| !gapinput@lad := LocalActionDiagramFromData(graph, vertex_labels, arc_labels);;|
  !gapprompt@gap>| !gapinput@View(lad);|
  <LocalActionDiagram with 1 vertex and 1 arc>
  !gapprompt@gap>| !gapinput@LocalActionDiagramIsRegularTreeAutomorphismGroup(lad);|
  true
  !gapprompt@gap>| !gapinput@View(lad);|
  <Aut(T5) (as a Local Action Diagram)>
\end{Verbatim}
 }

 }

   
\chapter{\textcolor{Chapter }{IO Operations and Visualisation}}\label{Chapter_IO_Operations_and_Visualisation}
\logpage{[ 3, 0, 0 ]}
\hyperdef{L}{X7E74125380A231F7}{}
{
  

 
\section{\textcolor{Chapter }{IO Operations}}\label{Chapter_IO_Operations_and_Visualisation_Section_IO_Operations}
\logpage{[ 3, 1, 0 ]}
\hyperdef{L}{X816FBAE6846CB59A}{}
{
  
\subsection{\textcolor{Chapter }{Writing To A String}}\label{AutoDoc_generated_group1}
\logpage{[ 3, 1, 1 ]}
\hyperdef{L}{X82FD5BB482591666}{}
{
\noindent\textcolor{FuncColor}{$\triangleright$\enspace\texttt{RSGraphToWritableString({\mdseries\slshape graph})\index{RSGraphToWritableString@\texttt{RSGraphToWritableString}!for IsRSGraph}
\label{RSGraphToWritableString:for IsRSGraph}
}\hfill{\scriptsize (operation)}}\\
\noindent\textcolor{FuncColor}{$\triangleright$\enspace\texttt{LocalActionDiagramToWritableString({\mdseries\slshape lad})\index{LocalActionDiagramToWritableString@\texttt{LocalActionDiagramToWritableString}!for IsLocalActionDiagram}
\label{LocalActionDiagramToWritableString:for IsLocalActionDiagram}
}\hfill{\scriptsize (operation)}}\\
\textbf{\indent Returns:\ }
A String representation of the object. 



 

 These two functions turn an RSGraph or local action diagram into a string
representation. The exact format of the representations are described in
(APPENDIX REFERENCE). These functions are not designed to store the objects in
the most space\texttt{\symbol{45}}efficient format. The formats are designed
so that the objects can be read into a GAP session with very little overhead
and so that the objects and known attributes can be recreated in a new GAP
session. 

 Note that these functions do not perform any disk IO operations themselves.
They return a string in the GAP session and it is up to the user to chose how
disk IO operations with the string are done. 

 For an RSGraph the data that is always stored is: 
\begin{itemize}
\item  the vertex ids, 
\item  the arc ids, 
\item  the adjacency list, and 
\item  the reverse map. 
\end{itemize}
 

 If the canonical labelling or automorphism group (REF FUNCTIONS) of the
RSGraph are known then they are also stored. 

 For a local action diagram that data that is always stored is: 
\begin{itemize}
\item  the diagrams RSGraph (and all data known about it), 
\item  the vertex labels, and 
\item  the arc labels. 
\end{itemize}
 

 The data stored if it is known is: 
\begin{itemize}
\item  the diagrams scopos, 
\item  the corresponding group type, 
\item  if the corresponding group is discrete, 
\item  if the corresponding group is uniscalar, 
\item  if the corresponding group is unimodular, and 
\item  the name of the corresponding group. 
\end{itemize}
 }

 
\begin{Verbatim}[commandchars=@CE,fontsize=\small,frame=single,label=Example]
  @gappromptCgap>E @gapinputCgraph := RSGraphByAdjacencyList([[1, 1]], ());;E
  @gappromptCgap>E @gapinputCRSGraphToWritableString(graph);E
  "1|1,1,1|P"
  
  @gappromptCgap>E @gapinputClad := LocalActionDiagramFromData(graph, [Group((1,2))], [[1,2]]);;E
  @gappromptCgap>E @gapinputCLocalActionDiagramGroupType(lad);;E
  @gappromptCgap>E @gapinputCLocalActionDiagramToWritableString(lad);E
  "1|1,1,1|P/1:P2,1:1,2|1:1,2|LocalActionDiagramGroupType!General|LocalActionDiagramScopos! "
\end{Verbatim}
 
\subsection{\textcolor{Chapter }{Reading From A String}}\label{AutoDoc_generated_group2}
\logpage{[ 3, 1, 2 ]}
\hyperdef{L}{X7D399A9A7C7BE4D7}{}
{
\noindent\textcolor{FuncColor}{$\triangleright$\enspace\texttt{RSGraphFromWritableString({\mdseries\slshape string})\index{RSGraphFromWritableString@\texttt{RSGraphFromWritableString}!for IsString}
\label{RSGraphFromWritableString:for IsString}
}\hfill{\scriptsize (operation)}}\\
\noindent\textcolor{FuncColor}{$\triangleright$\enspace\texttt{LocalActionDiagramFromWritableString({\mdseries\slshape string[, return{\textunderscore}graph]})\index{LocalActionDiagramFromWritableString@\texttt{Local}\-\texttt{Action}\-\texttt{Diagram}\-\texttt{From}\-\texttt{Writable}\-\texttt{String}!for IsString, IsBool}
\label{LocalActionDiagramFromWritableString:for IsString, IsBool}
}\hfill{\scriptsize (operation)}}\\
\textbf{\indent Returns:\ }
The object(s) the string represents. 



 

 These function take the output from the \texttt{RSGraphToWritableString} and \texttt{LocalActionDiagramToWritableString"} methods as input and returns the objects those strings represents. While it is
possible to edit these strings outside of GAP these functions do not perform
error checking and so this can potentially lead to invalid RSGraph or local
action diagram objects. If the optional \mbox{\texttt{\mdseries\slshape return{\textunderscore}graph}} parameter is set to true then the \texttt{LocalActionDiagramFromWritableString} function will return a list whose first entry is the local action diagram
object and second intro is the RSGraph for this local action diagram. 

 Any optional attributes known about the objects at the time of writing will be
stored in the objects returned by these functions without needing to be
recomputed. Note that the (REF canon and aut functions) attributes can be
stored in the objects created by these functions without needing the \textsf{Digraphs} to be loaded even though computing these attributes required the \textsf{Digraphs} to be loaded. }

 
\begin{Verbatim}[commandchars=@BC,fontsize=\small,frame=single,label=Example]
  @gappromptBgap>C @gapinputBgraph_string := "1|1,1,1|P";;C
  @gappromptBgap>C @gapinputBgraph := RSGraphFromWritableString(graph_string);;C
  @gappromptBgap>C @gapinputBPrint(graph);C
  Vertices = { 1 }
  Arcs = {
      1 = ( origin = 1, terminus = 1, inverse = 1 )
  }
  Reverse Map = ()
  
  @gappromptBgap>C @gapinputBlad_string := "1|1,1,1|P/1:P2,1:1,2|1:1,2|LocalActionDiagramGroupType!General|LocalActionDiagramScopos! ";;C
  @gappromptBgap>C @gapinputBlad := LocalActionDiagramFromWritableString(lad_string);C
  @gappromptBgap>C @gapinputBPrint(lad);C
  Vertices = { 1 }
  Arcs = {
      1 = ( origin = 1, terminus = 1, inverse = 1 )
  }
  Reverse Map = ()
  Vertex Labels = {
      1 = Group( [ (1,2) ] )
  }
  Arc Labels = {
      1 = [ 1, 2 ]
  }
  @gappromptBgap>C @gapinputBPrint("LocalActionDiagramScopos" in KnownAttributesOfObject(lad));C
  true
  @gappromptBgap>C @gapinputBPrint("LocalActionDiagramIsDiscrete" in KnownAttributesOfObject(lad));C
  false
\end{Verbatim}
 
\subsection{\textcolor{Chapter }{IO Pickling/Unpickling}}\logpage{[ 3, 1, 3 ]}
\hyperdef{L}{X7E57D3CD7869D64C}{}
{
\noindent\textcolor{FuncColor}{$\triangleright$\enspace\texttt{IO{\textunderscore}Pickle({\mdseries\slshape [file, ]graph})\index{IO{\textunderscore}Pickle@\texttt{IO{\textunderscore}Pickle}!for [an IO File and] an RSGraph}
\label{IOuScorePickle:for [an IO File and] an RSGraph}
}\hfill{\scriptsize (operation)}}\\
\noindent\textcolor{FuncColor}{$\triangleright$\enspace\texttt{IO{\textunderscore}Pickle({\mdseries\slshape [file, ]lad})\index{IO{\textunderscore}Pickle@\texttt{IO{\textunderscore}Pickle}!for [an IO File and] a Local Action Diagram}
\label{IOuScorePickle:for [an IO File and] a Local Action Diagram}
}\hfill{\scriptsize (operation)}}\\
\textbf{\indent Returns:\ }
\texttt{IO{\textunderscore}OK} or the pickle string if the pickling was successful or \texttt{IO{\textunderscore}ERROR} if there was an error.

\noindent\textcolor{FuncColor}{$\triangleright$\enspace\texttt{IO{\textunderscore}Unpickle({\mdseries\slshape file})\index{IO{\textunderscore}Unpickle@\texttt{IO{\textunderscore}Unpickle}!for an IO File or String}
\label{IOuScoreUnpickle:for an IO File or String}
}\hfill{\scriptsize (operation)}}\\
\textbf{\indent Returns:\ }
An RSGraph or Local Action Diagram object if successful or \texttt{IO{\textunderscore}ERROR} if there was an error.



 The IO package supports serialising data through the \texttt{IO{\textunderscore}Pickle} function. Unlike the \texttt{RSGraphToWritableString} and \texttt{LocalActionDiagramToWritableString} functions these functions support recreating an arbitrary structure within the
memory of a GAP session (assuming that the pickling functions have been
defined for each object). In particular, it can easily recreate data
structures storing objects. 

 For example, if \texttt{l} is a list of RSGraphs then you can write this list to the disk with a single
call to \texttt{IO{\textunderscore}Pickle}. This generality comes at the cost of performance. For this reason, if you
have a large collection of RSGraph and local action diagram objects that you
need writing to the disk we recommend using the \texttt{RSGraphToWritableString} and \texttt{LocalActionDiagramToWritableString} functions. 

 In the one argument version for \texttt{IO{\textunderscore}Pickle}, the pickle string for the object is returned if the pickling was successful.
In the two argument version, an IO file opened to write mode is supplied and
the pickle string is written to this file. 

 The \texttt{IO{\textunderscore}Unpickle} function accepts either the pickle string as input or an IO file opened to
read mode which stores pickled objects. It recreates the objects in memory
exactly. 

 As a note for developers, the RSGraph and local action diagram objects are
pickled with the "magic values" \texttt{RSGO} and \texttt{LADO} respectively. }

 
\begin{Verbatim}[commandchars=!@|,fontsize=\small,frame=single,label=Example]
  !gapprompt@gap>| !gapinput@graph := RSGraphByAdjacencyList([[1, 1]], ());;|
  !gapprompt@gap>| !gapinput@lad := LocalActionDiagramFromData(graph, [Group((1,2))], [[1,2]]);;|
  !gapprompt@gap>| !gapinput@list := [graph, lad];|
  [ <RSGraph with 1 vertex and 1 arc>, <LocalActionDiagram with 1 vertex and 1 arc> ]
  !gapprompt@gap>| !gapinput@file := IO_File("test.pickle", "w");;|
  !gapprompt@gap>| !gapinput@IO_Pickle(file, list);|
  IO_OK
  !gapprompt@gap>| !gapinput@IO_Close(file);;|
  !gapprompt@gap>| !gapinput@file := IO_File("test.pickle", "r");;|
  !gapprompt@gap>| !gapinput@IO_Unpickle(file);|
  [ <RSGraph with 1 vertex and 1 arc>, <LocalActionDiagram with 1 vertex and 1 arc> ]
\end{Verbatim}
 }

 }

   
\chapter{\textcolor{Chapter }{Isomorphisms and Automorphisms}}\label{Chapter_Isomorphisms_and_Automorphisms}
\logpage{[ 4, 0, 0 ]}
\hyperdef{L}{X81F3496578EAA74E}{}
{
  

 }

   
\chapter{\textcolor{Chapter }{Enumeration}}\label{Chapter_Enumeration}
\logpage{[ 5, 0, 0 ]}
\hyperdef{L}{X872DDB2582015ED5}{}
{
  

 }

 \def\indexname{Index\logpage{[ "Ind", 0, 0 ]}
\hyperdef{L}{X83A0356F839C696F}{}
}

\cleardoublepage
\phantomsection
\addcontentsline{toc}{chapter}{Index}


\printindex

\immediate\write\pagenrlog{["Ind", 0, 0], \arabic{page},}
\newpage
\immediate\write\pagenrlog{["End"], \arabic{page}];}
\immediate\closeout\pagenrlog
\end{document}
